% =====================================================================
% 第一讲 (Exposé I) 平展态射 (Morphismes étales)
% A. Grothendieck
% 原书页码: pp. 1-24
% =====================================================================
\chapter{平展态射}

\section{微分演算的概念}

设 $X$ 是 $Y$ 上的\emph{预概形} (prescheme)\index[term]{yugaixing@预概形}，设 $\Delta_{X/Y}$ 或 $\Delta$ 是对角态射 (diagonal morphism) $X \to X \times_Y X$. 它是\emph{浸入} (immersion) 态射，因而是从 $X$ 到 $X \times_Y X$ 的某个开集 $V$ 中的闭浸入态射. 设 $\mathcal{I}_X$ 是 $V$ 中对应于该对角线的闭子预概形的理想 (N.B. 若想作内蕴的处理，而不假设 $X$ 在 $Y$ 上分离 (separated) —— 这个假设会是荒谬的 (canularesque) —— 就应当考虑 $\mathcal{O}_{X \times X}$ 在 $X$ 中的集合论原像，并用 $\mathcal{I}_X$ 表示后者中的增广理想 (idéal d'augmentation)...). \emph{层} (sheaf) $\mathcal{I}_X / \mathcal{I}_X^2$ 可以看作 $X$ 上的\emph{拟凝聚层} (quasi-coherent sheaf)，记为 $\varOmega_{X/Y}^{1}$. 若 $X \to Y$ 是有限型 (of finite type) 的，则它是有限型的. 它对基扩张 (base extension) $Y' \to Y$ 表现良好. 还引入层 $\mathcal{O}_{X \times_Y X} / \mathcal{I}_X^{n+1} = \mathcal{P}_{X/Y}^{n}$，它们是 $X$ 上的环层，使 $X$ 成为一个预概形，可记为 $\Delta_{X/Y}^{n}$ 并称为 $X/Y$ 的第 $n$ 个\emph{无穷小邻域} (infinitesimal neighborhood)\index[term]{wuqiongxiaolinyu@无穷小邻域}. 关于它的这套推论 (sorites) 完全平凡，尽管相当长\footnote{参见 EGA IV 16.3.}；谨慎起见，最好还是等到我们对\emph{光滑态射} (smooth morphism) 有可用的结论时再谈它.

\section{拟有限态射}

\begin{proposition}\label{prop:I.2.1}
设 $A \to B$ 是局部同态 (local homomorphism) (N.B. 现在环都是诺特 (noetherian) 的)，$m$ 是 $A$ 的极大理想 (maximal ideal). 下列条件等价：
\begin{enumerate}[(i)]
\item $B/mB$ 在 $k = A/m$ 上是有限维的；
\item $mB$ 是定义理想 (ideal of definition)，且 $B/\mathfrak{r}(B) = k(B)$ 是 $k = k(A)$ 的扩张；
\item $B$ 的完备化 (completion) $\widehat{B}$ 在 $A$ 的完备化 $\widehat{A}$ 上有限.
\end{enumerate}
\end{proposition}

这时称 $B$ 在 $A$ 上拟有限 (quasi-finite). 态射 $f \colon X \to Y$ 称为在 $x$ 处拟有限 (或 $Y$-预概形\index[term]{niyouxiantaishe@拟有限态射} $f$ 称为在 $x$ 处拟有限)，若 $\mathcal{O}_x$ 在 $\mathcal{O}_{f(x)}$ 上拟有限. 这也等价于说 $x$ 在其纤维 (fiber) $f^{-1}(x)$ 中是孤立的 (isolated). 态射称为拟有限的，若它在每一点处都是拟有限的\footnote{在 EGA II 6.2.3 中还进一步假设 $f$ 是有限型的.}.

\begin{corollary}\label{cor:I.2.2}
若 $A$ 完备，则拟有限等价于有限.
\end{corollary}

我们本可以给出\emph{拟有限态射}的通常那套推论 (i)，(ii)，(iii)，(iv)，(v)，但在这里似乎没有必要.

\section{非分歧态射}

\begin{proposition}\label{prop:I.3.1}
设 $f \colon X \to Y$ 是有限型态射，$x \in X$，$y = f(x)$. 下列条件等价：
\begin{enumerate}[(i)]
\item $\mathcal{O}_x / m_y \mathcal{O}_x$ 是 $k(y)$ 的有限可分扩张 (finite separable extension)；
\item $\varOmega_{X/Y}^{1}$ 在 $x$ 处为零；
\item 对角态射 $\Delta_{X/Y}$ 在 $x$ 的邻域中是开浸入.
\end{enumerate}
\end{proposition}

\begin{proof}
对于蕴含关系 (i)$\Rightarrow$(ii)，由 Nakayama 引理立即归结到 $Y = \mathrm{Spec}(k)$，$X = \mathrm{Spec}(k')$ 的情形，在那里这是众所周知的，而且从可分的定义看甚至是平凡的；(ii)$\Rightarrow$(iii) 依据一个用 Krull 定理的开浸入的巧妙而容易的刻画；(iii)$\Rightarrow$(i) 是因为我们仍归结到 $Y = \mathrm{Spec}(k)$ 且对角态射处处是开浸入的情形. 于是必须证明 $X$ 是 $k$ 上的有限可分环代数，为此归结到 $k$ 代数闭 (algebraically closed) 的情形. 但此时 $X$ 的每个闭点都是孤立的 (因为它等于由 $x$ 定义的态射 $X \to X \times_k X$ 之下对角线的原像)，故 $X$ 有限. 然后可以假设 $X$ 化为一个点，其环为 $A$，从而 $A \otimes_k A \to A$ 是同构，故 $A = k$，证毕.
\end{proof}

\begin{definition}\label{def:I.3.2}
\begin{enumerate}[a)]
\item 这时称 $f$ 在 $x$ 处是非分歧的 (unramified)\index[term]{feifenqitaishe@非分歧态射} (unramified)，或称 $X$ 在 $Y$ 上于 $x$ 处非分歧，或非分歧.
\item 设 $A \to B$ 是局部同态，称它是非分歧的，或称 $B$ 是 $A$ 上的非分歧局部代数\index[term]{feifenqidaishu@非分歧代数} (net local algebra)，若 $B/\mathfrak{r}(A)B$ 是 $A/\mathfrak{r}(A)$ 的有限可分扩张，即若 $\mathfrak{r}(A)B = \mathfrak{r}(B)$ 且 $k(B)$ 是 $k(A)$ 的可分扩张\footnote{参见 \ref{rem:III.1.2} 中的更正 (retraction).}.
\end{enumerate}
\end{definition}

\begin{remarkstar}
$B$ 在 $A$ 上非分歧这一事实从 $A$ 与 $B$ 的完备化上即可看出. 非分歧蕴含拟有限.
\end{remarkstar}

\begin{corollary}\label{cor:I.3.3}
$f$ 为非分歧的那些点构成的集合是开的.
\end{corollary}

\begin{corollary}\label{cor:I.3.4}
设 $X'$，$X$ 是 $Y$ 上的两个有限型预概形，$g \colon X' \to X$ 是 $Y$-态射. 若 $X$ 在 $Y$ 上非分歧，则图像态射 (graph morphism) $\Gamma_g \colon X' \to X \times_Y X$ 是开浸入.
\end{corollary}

\begin{proof}
事实上，它就是对角态射 $X \to X \times_Y X$ 在
\[
g \times_Y \mathrm{id}_{X'} \colon X' \times_Y X \longrightarrow X \times_Y X
\]
之下的原像.
\end{proof}

还可以引入 $\varOmega_{X/Y}^{1}$ 的零化理想 $\mathfrak{d}_{X/Y}$，称为 $X/Y$ 的差理想 (different ideal)\index[term]{panbielixiang@判别理想} 或微分歧理想；它定义 $X$ 的一个闭子预概形，从集合观点看，它就是 $X/Y$ 分歧 (ramified)，即非非分歧的那些点构成的集合.

\begin{proposition}\label{prop:I.3.5}
\begin{enumerate}[(i)]
\item 浸入是非分歧的.\index[term]{jinru@浸入}
\item 两个\emph{非分歧态射}的复合是非分歧的.
\item 非分歧态射的基扩张 (base extension) 是非分歧的.
\end{enumerate}
\end{proposition}

\begin{proof}
可以从 (ii) 或 (iii) 直接看出 (第二个在我看来更有趣). 当然也可以把它们细化成逐点的陈述；这只是表面上更一般 (除定义 b) 的框架外)，而且令人厌烦.
\end{proof}

照例得到推论：

\begin{corollary}\label{cor:I.3.6}
\begin{enumerate}[(iv)]
\item 两个非分歧态射的笛卡儿积 (Cartesian product) 是非分歧的.
\item 若 $g \circ f$ 非分歧，则 $f$ 非分歧.
\item 若 $f$ 非分歧，则 $f_{\mathrm{r\acute{e}d}}$ 非分歧.
\end{enumerate}
\end{corollary}

\begin{proposition}\label{prop:I.3.7}
设 $A \to B$ 是局部同态，假设剩余扩张 (residue field extension) $k(B)/k(A)$ 平凡，或 $k(A)$ 代数闭. 为使 $B/A$ 非分歧，必要且充分条件是 $\widehat{B}$ (作为 $A$-代数) 是 $\widehat{A}$ 的商.
\end{proposition}

\begin{remarkstar}
在不假设剩余扩张平凡的情形，可以对 $A$ 作一个适当的有限平坦扩张，以破坏所说的扩张，从而归结到这一情形.

举例：设 $A$ 是曲线的一个普通二重点 (ordinary double point) 的局部环，$B$ 是其正规化上一个点的局部环：那么 $A \subset B$，$B$ 在 $A$ 上非分歧且剩余扩张平凡，而 $\widehat{A} \to \widehat{B}$ 是满射但非单射. 因此我们要强化非分歧性的概念.
\end{remarkstar}

\section{平展态射。平展覆盖}

我们将承认关于平坦 (flat) 态射所需的一切；这些事实将在稍后予以证明，如果有必要的话\footnote{这里原书注 (4) 缺失.}.

\begin{definition}\label{def:I.4.1}
\begin{enumerate}[a)]
\item 设 $f \colon X \to Y$ 是有限型态射. 称 $f$ 在 $x$ 处平展 (etale)，若 $f$ 在 $x$ 处平坦且在 $x$ 处非分歧. 称 $f$ 平展，若它在所有点处平展. 称 $X$ 在 $Y$ 上于 $x$ 处平展，或它是 $x$ 处的平展 $Y$-预概形，等等.
\item 设 $f \colon A \to B$ 是局部同态. 称 $f$ 平展，或称 $B$ 在 $A$ 上平展，若 $B$ 在 $A$ 上平坦且非分歧\footnote{原书注 (5) 缺失.}.
\end{enumerate}
\end{definition}

\begin{proposition}\label{prop:I.4.2}
为使 $B/A$ 平展，必要且充分条件是 $\widehat{B}/\widehat{A}$ 平展.
\end{proposition}

\begin{proof}
事实上，对``非分歧''和``平坦''分别成立即可.
\end{proof}

\begin{corollary}\label{cor:I.4.3}
设 $f \colon X \to Y$ 是有限型的，$x \in X$. $f$ 在 $x$ 处平展这一事实只依赖于局部同态 $\mathcal{O}_{f(x)} \to \mathcal{O}_x$，甚至只依赖于对完备化所作的相应同态.
\end{corollary}

\begin{corollary}\label{cor:I.4.4}
假设剩余扩张 $k(A) \to k(B)$ 平凡，或 $k(A)$ 代数闭. 则 $B/A$ 平展当且仅当 $\widehat{A} \to \widehat{B}$ 是同构.
\end{corollary}

\begin{proof}
把平坦性与 \ref{prop:I.3.7} 结合起来即得.
\end{proof}

\begin{proposition}\label{prop:I.4.5}
设 $f \colon X \to Y$ 是有限型态射. 则 $f$ 为平展的那些点构成的集合是开的.
\end{proposition}

\begin{proof}
事实上，对``非分歧''和``平坦''分别成立即可.
\end{proof}

本命题说明，在研究处处平展的有限型态射时，可以放弃``逐点''的陈述.

\begin{proposition}\label{prop:I.4.6}
\begin{enumerate}[(i)]
\item 开浸入是平展的.
\item 两个平展态射的复合是平展的.
\item 基扩张 (base extension).
\end{enumerate}
\end{proposition}

\begin{proof}
事实上，(i) 是平凡的，而对 (ii) 和 (iii)，只需注意这对``非分歧''和``平坦''成立. 老实说，对局部同态 (无有限性条件) 也有相应的陈述，它们在各种情形都应当进入大全 (multiplodoque)，首先是非分歧的情形.
\end{proof}

\begin{corollary}\label{cor:I.4.7}
两个平展态射的笛卡儿积也是平展的.
\end{corollary}

\begin{corollary}\label{cor:I.4.8}
设 $X$ 和 $X'$ 在 $Y$ 上有限型，$g \colon X \to X'$ 是 $Y$-态射. 若 $X'$ 在 $Y$ 上非分歧且 $X$ 在 $Y$ 上平展，则 $g$ 平展.
\end{corollary}

\begin{proof}
事实上，$g$ 是图像态射 $\Gamma_g \colon X \to X \times_Y X'$ (由 \ref{cor:I.3.4} 它是开浸入) 与投影态射的复合，而后者平展，因为它由平展态射 $X \to Y$ 经``基扩张'' $X' \to Y$ 得到.
\end{proof}

\begin{definition}\label{def:I.4.9}
称 $Y$ 的平展覆盖 (etale covering) (相应地：非分歧覆盖) 为 $Y$ 上的一个 $Y$-概形 $X$，它在 $Y$ 上有限且平展 (相应地：非分歧).
\end{definition}

第一个条件意味着 $X$ 由 $Y$ 上的一个凝聚 (coherent) 代数层 $\mathcal{B}$ 定义. 第二个条件则意味着 $\mathcal{B}$ 在 $Y$ 上局部自由 (locally free) (相应地：什么也不意味着)，并且此外对每个 $y \in Y$，纤维 (fiber) $\mathcal{B}(y) = \mathcal{B}_y \otimes_{\mathcal{O}_y} k(y)$ 是 $k(y)$ 上的可分代数 (= 有限个有限可分扩张的直积).

\begin{proposition}\label{prop:I.4.10}
设 $X$ 是 $Y$ 的 $n$ 次平坦覆盖 (这个术语的定义本应放在 \ref{def:I.4.9} 中)，由凝聚局部自由代数层 $\mathcal{B}$ 定义. 以众所周知的方式定义迹同态 (trace homomorphism) $\mathcal{B} \to \mathcal{A}$ (它是 $\mathcal{A}$-模同态，其中 $\mathcal{A} = \mathcal{O}_Y$). 为使 $X$ 平展，必要且充分条件是相应的双线性型 (bilinear form) $\mathrm{tr}_{\mathcal{B}/\mathcal{A}}\, xy$ 定义 $\mathcal{B}$ 到其对偶 $\check{\mathcal{B}}$ 的同构，或等价地，判别式 (discriminant) 截面
\[
d_{X/Y} = d_{\mathcal{B}/\mathcal{A}} \in \Gamma\bigl(Y, \textstyle{\bigwedge^{n}} \check{\mathcal{B}} \otimes_{\mathcal{A}} \textstyle{\bigwedge^{n}} \check{\mathcal{B}}\bigr)
\]
可逆，或最后，由该截面定义的\emph{判别理想} (discriminant ideal) 是单位理想.
\end{proposition}

\begin{proof}
事实上，归结到 $Y = \mathrm{Spec}(k)$ 的情形，此时这是众所周知的可分性判则，通过 $k$ 的代数闭包 (algebraic closure) 即可平凡地看出.
\end{proof}

\begin{remarkstar}
下面，当不预先假设 $X$ 在 $Y$ 上平坦而只作正规性 (normality) 假设时，我们有一个不那么平凡的陈述.
\end{remarkstar}

\section{平展态射的基本性质}

\begin{theorem}\label{thm:I.5.1}
设 $f \colon X \to Y$ 是有限型态射. 为使 $f$ 是开浸入，必要且充分条件是它是平展且根式 (radicial) 的态射.
\end{theorem}

\begin{proof}
回忆根式意为：单射，且剩余扩张是根式的 (还可以回忆这等价于：该态射对任何基扩张仍是单射). 必要性是平凡的，剩下充分性. 我们给出两个不同的证明，第一个更短，第二个更初等.
\end{proof}

\begin{enumerate}[1)]
\item 平坦态射是开的，因此可以假设 (把 $Y$ 换成 $f(X)$) $f$ 是满的同胚 (homeomorphism). 对任何基扩张，$f$ 仍将平坦、根式、满射，因而仍是同胚，从而是闭的. 故 $f$ \emph{紧合} (proper). 因此 $f$ 有限 (参考文献：Chevalley 定理)，由凝聚代数层 $\mathcal{B}$ 定义. $\mathcal{B}$ 局部自由，且由假设它在处处秩 (rank) 为 1，故 $X = Y$，证毕.
\item 可以假设 $Y$ 和 $X$ 仿射 (affine). 我们更容易归结到证明：若 $Y = \mathrm{Spec}(A)$，$A$ 局部，且 $f^{-1}(y)$ 非空 ($y$ 是 $Y$ 的闭点)，则 $X = Y$ (事实上，这将蕴含每个 $y \in f(X)$ 有一个开邻域 $U$ 使得 $X|U = U$). 设 $X = \mathrm{Spec}(B)$，我们要证明 $A = B$. 为此归结到把 $A$ 换成 $\widehat{A}$，$B$ 换成 $B \otimes_A \widehat{A}$ 的相应断言 (考虑到 $\widehat{A}$ 在 $A$ 上忠实平坦 (faithfully flat)). 故可假设 $A$ 完备. 设 $x$ 是 $y$ 上的点，由推论 \ref{cor:I.2.2}，$\mathcal{O}_x$ 在 $A$ 上有限，从而 (在 $A$ 上平坦且根式) 与 $A$ 相同. 于是 $X = Y \amalg X'$ (不交并). 由于 $X$ 在 $Y$ 上根式，$X'$ 为空. 证毕.
\end{enumerate}

\begin{corollary}\label{cor:I.5.2}
设 $f \colon X \to Y$ 是闭浸入且平展的态射. 若 $X$ 连通，则 $f$ 是 $X$ 到 $Y$ 的一个连通分支上的同构.
\end{corollary}

\begin{proof}
事实上，$f$ 也是开浸入. 由此得出：
\end{proof}

\begin{corollary}\label{cor:I.5.3}
设 $X$ 是 $Y$ 上的非分歧概形，$Y$ 连通. 则 $X$ 在 $Y$ 上的每个截面 (section) 都是 $Y$ 到 $X$ 的一个连通分支上的同构. 因此在所有这些截面的集合与 $X$ 的满足``投影 $X_i \to Y$ 是同构''(或由 \ref{thm:I.5.1} 等价地：满射且根式) 的连通分支 $X_i$ 的集合之间存在一一对应. 特别地，一个截面由其在一个点的值即可确定.
\end{corollary}

\begin{proof}
只有第一个断言需要证明；由 \ref{cor:I.5.2}，只需注意一个截面是闭浸入 (因为 $X$ 在 $Y$ 上分离) 且由 \ref{cor:I.4.8} 它是平展的.
\end{proof}

\begin{corollary}\label{cor:I.5.4}
设 $X$ 和 $Y$ 是 $S$ 上的两个预概形，$X$ 在 $S$ 上非分歧且分离，$Y$ 连通. 设 $f$，$g$ 是 $Y$ 到 $X$ 的两个 $S$-态射，$y$ 是 $Y$ 的一个点，假设 $f(y) = g(y) = x$，且由 $f$ 和 $g$ 定义的剩余同态 $k(x) \to k(y)$ 相同 (<<$f$ 和 $g$ 在 $y$ 处几何上重合>>). 则 $f$ 和 $g$ 相同.
\end{corollary}

\begin{proof}
由 \ref{cor:I.5.3} 归结到 $Y = S$ 的情形，把 $X$ 换成 $X \times_S Y$.
\end{proof}

下面是 \ref{cor:I.5.3} 的一个特别重要的变体.

\begin{theorem}\label{thm:I.5.5}
设 $S$ 是预概形，$X$ 和 $Y$ 是两个 $S$-预概形，$S_0$ 是 $S$ 的一个与 $S$ 有相同底空间的闭子预概形，$X_0 = X \times_S S_0$ 和 $Y_0 = Y \times_S S_0$ 是 $X$ 和 $Y$ 在 $S_0$ 上的``限制''. 假设 $X$ 在 $S$ 上平展. 则自然映射
\[
\operatorname{Hom}_{S}(Y, X) \longrightarrow \operatorname{Hom}_{S_0}(Y_0, X_0)
\]
是双射.
\end{theorem}

\begin{proof}
仍归结到 $Y = S$ 的情形，于是由 \ref{cor:I.5.3} 中给出的 $X/Y$ 的截面的``拓扑''描述即得.
\end{proof}

\begin{scholie}
这个结果包含态射的存在性与唯一性断言. 它还可以表述为 (当 $X$ 和 $Y$ 都在 $S$ 上平展时)：从平展 $S$-概形的范畴到平展 $S_0$-概形的范畴的函子 $X \mapsto X_0$ 是全忠实 (fully faithful) 的，即建立起第一个范畴与第二个范畴的一个全子范畴 (full subcategory) 之间的等价. 我们将在下面看到，它甚至是第一个与第二个范畴之间的等价 (这将是一个平展 $S$-概形的存在性定理).
\end{scholie}

下面是 \ref{thm:I.5.5} 的表面上更一般的形式，它常常是方便的：

\begin{corollary}[«提升延拓定理» (Theorem of extension of liftings)]\label{cor:I.5.6}
考虑态射的交换图
\[
\begin{tikzcd}[row sep=large, column sep=large]
X & Y_0 \arrow[l] \arrow[d, hook] \\
S & Y \arrow[l] \arrow[u, dashed, "\exists!"]
\end{tikzcd}
\]
其中 $X \to S$ 平展，$Y_0 \to Y$ 是双射闭浸入. 则可以找到唯一的态射 $Y \to X$，使相应的两个三角形交换.
\end{corollary}

\begin{proof}
事实上，把 $S$ 换成 $Y$，$X$ 换成 $X \times_S Y$，归结到 $Y = S$ 的情形，此时它是 \ref{thm:I.5.5} 对 $Y = S$ 的特殊情形.
\end{proof}

再指出 \ref{thm:I.5.1} 的如下直接推论 (我们当初未把它列为推论，是为了不打断 \ref{thm:I.5.1} 之后所发展的思路)：

\begin{proposition}\label{prop:I.5.7}
设 $X$，$X'$ 是两个在 $Y$ 上有限型且平坦的预概形，$g \colon X \to X'$ 是 $Y$-态射. 为使 $g$ 是开浸入 (相应地：同构)，必要且充分条件是对每个 $y \in Y$，在纤维上诱导的态射
\[
g \otimes_Y k(y) \colon X \otimes_Y k(y) \longrightarrow X' \otimes_Y k(y)
\]
是开浸入 (相应地：同构).
\end{proposition}

\begin{proof}
只需证明充分性；由于这对满射概念成立，归结到开浸入的情形. 由 \ref{thm:I.5.1}，必须验证 $g$ 根式 (这是平凡的) 且平展 (这由下面的推论 \ref{cor:I.5.9} 得出).
\end{proof}

\begin{corollary}\label{cor:I.5.8}
(本应移到第 3 节) 设 $X$ 和 $X'$ 是两个 $Y$-预概形，$g \colon X \to X'$ 是 $Y$-态射，$x$ 是 $X$ 的点，$y$ 是它在 $Y$ 上的投影. 为使 $g$ 在 $x$ 处拟有限 (相应地：非分歧)，必要且充分条件是 $g \otimes_Y k(y)$ 亦然.
\end{corollary}

\begin{proof}
事实上，为了确认在 $x$ 处确有拟有限 (相应地：非分歧) 态射而需要考察的 $k(g(x))$ 上的两个代数，对 $g$ 和 $g \otimes_Y k(y)$ 是相同的.
\end{proof}

\begin{corollary}\label{cor:I.5.9}
沿用 \ref{cor:I.5.8} 的记号，假设 $X$ 和 $X'$ 在 $Y$ 上平坦且有限型. 为使 $g$ 在 $x$ 处平坦 (相应地：平展)，必要且充分条件是 $g \otimes_Y k(y)$ 平坦 (相应地：平展).
\end{corollary}

\begin{proof}
对于``平坦''，这个陈述只是备忘，它是平坦性的基本判则之一\footnote{参见 EGA IV 11.8.2 (原书注 (6)).}. 对于平展，结合 \ref{cor:I.5.8} 即得.
\end{proof}

\section{完备局部环的平展扩张的应用}

本节是关于形式预概形 (formal prescheme) 的结果的一个特殊情形，那些结果应当进入大全 (multiplodoque). 不过在这里我们更省事，即不需要 §7 中平展态射的显式局部刻画 (它使用主定理 (Main Theorem)). 这也许是保留本节 (即使在大全中也保留) 在此位置的充分理由.

\begin{theorem}\label{thm:I.6.1}
设 $A$ 是完备局部环 (当然是诺特的)，剩余域为 $k$. 对每个 $A$-代数 $B$，设 $R(B) = B \otimes_A k$ 看作 $k$-代数，它函子性地依赖于 $B$. 则 $R$ 建立 $A$ 上有限且平展的 $A$-代数的范畴与 $k$ 上有限阶可分代数的范畴之间的等价.
\end{theorem}

\begin{proof}
首先，所述函子全忠实，这由下述更一般的结论得出：
\end{proof}

\begin{corollary}\label{cor:I.6.2}
设 $B$，$B'$ 是两个在 $A$ 上有限的 $A$-代数. 若 $B$ 在 $A$ 上平展，则典范映射
\[
\operatorname{Hom}_{A\text{-alg}}(B, B') \longrightarrow \operatorname{Hom}_{k\text{-alg}}\bigl(R(B), R(B')\bigr)
\]
是双射.
\end{corollary}

\begin{proof}
归结到 $A$ 为 Artin 环的情形 (把 $A$ 换成 $A/\mathfrak{m}^n$)，此时它是 \ref{thm:I.5.5} 的特殊情形.
\end{proof}

剩下要证明：对每个有限可分 $k$-代数 (为什么不说：平展，这更短) $L$，存在平展于 $A$ 上的 $B$ 使得 $R(B)$ 同构于 $L$. 可以假设 $L$ 是 $k$ 的可分扩张，作为这样的扩张它有一个生成元 (generator) $x$，即同构于一个代数 $k[t]/Fk[t]$，其中 $F \in k[t]$ 是首一多项式 (monic polynomial). 把 $F$ 提升为首一多项式 $F_1 \in A[t]$，并取 $B = A[t]/F_1 A[t]$.

\section{非分歧态射与平展态射的局部构造}

\begin{proposition}\label{prop:I.7.1}
设 $A$ 是诺特环，$B$ 是 $A$ 上的有限代数，$u$ 是 $B$ 在 $A$ 上的生成元，$F \in A[t]$ 满足 $F(u) = 0$ (不假设 $F$ 首一)，$u' = F'(u)$ (其中 $F'$ 是导多项式 (derived polynomial))，$q$ 是 $B$ 的一个不含 $u'$ 的素理想，$p$ 是它在 $A$ 上的迹. 则 $B_{\mathfrak{q}}$ 在 $A_{\mathfrak{p}}$ 上非分歧.
\end{proposition}

\begin{proof}
换言之，令 $Y = \mathrm{Spec}(A)$，$X = \mathrm{Spec}(B)$，$X_{u'} = \mathrm{Spec}(B_{u'})$，则 $X_{u'}$ 在 $Y$ 上非分歧. 本陈述由下面更精确的结论得出：
\end{proof}

\begin{corollary}\label{cor:I.7.2}
$B/A$ 的差理想包含 $u'B$，并且若自然同态 $A[t]/FA[t] \to B$ (把 $t$ 映到 $u$) 是同构，则与之相等.
\end{corollary}

\begin{proof}
设 $J$ 是同态 $C = A[t] \to B$ 的核，它包含 $FA[t]$，并且在 \ref{cor:I.7.2} 所考虑的第二种情形中与之相等. 由于它是满射，$\varOmega_{B/A}^{1}$ 等同于 $\varOmega_{C/A}^{1}$ 除以由 $J\varOmega_{C/A}^{1}$ 和 $d(J)$ 生成的子模 (本应在 §1 中详细说明同态 $d$ 的定义，以及多项式代数的 $\varOmega^{1}$ 的计算). 借助基 $dt$ 把 $\varOmega_{C/A}^{1}$ 等同于 $C$，得到 $B / B \cdot J'$，故差理想由 $J$ 中诸 $G$ 的导数在 $B$ 中的像的集合 $J'$ 生成 (只需取生成 $J$ 的 $G$). 由于 $F \in J$ (相应地：$F$ 生成 $J$)，证毕. (N.B. 应当把 \ref{cor:I.7.2} 作为命题、把 \ref{prop:I.7.1} 作为推论). 得到：
\end{proof}

\begin{corollary}\label{cor:I.7.3}
在 \ref{prop:I.7.1} 的条件下，假设 $F$ 首一且 $A[t]/FA[t] \to B$ 是同构，为使 $B_{\mathfrak{q}}$ 在 $A_{\mathfrak{p}}$ 上平展，必要且充分条件是 $q$ 不含 $u'$.
\end{corollary}

\begin{proof}
事实上，由于 $B$ 在 $A$ 上平坦，平展等价于非分歧，并可应用 \ref{cor:I.7.2}.
\end{proof}

\begin{corollary}\label{cor:I.7.4}
在 \ref{cor:I.7.3} 的条件下，为使 $B$ 在 $A$ 上平展，必要且充分条件是 $u'$ 可逆，或等价地，$F$，$F'$ 在 $A[t]$ 中生成的理想是单位理想.
\end{corollary}

\begin{proof}
最后一个判则由第一个判则和 Nakayama 引理 (在 $B$ 中) 得出.
\end{proof}

具有推论 \ref{cor:I.7.4} 所述性质的首一多项式 $F \in A[t]$ 称为\emph{可分多项式} (separable polynomial)\index[term]{kefenduoxiangshi@可分多项式} (若 $F$ 非首一，至少应要求其首项系数可逆；当 $A$ 是域时，就得到通常的定义).

\begin{corollary}\label{cor:I.7.5}
设 $B$ 是局部环 $A$ 上的有限代数. 假设 $k(A)$ 无限，或 $B$ 局部. 设 $n$ 是 $L = B \otimes_A k(A)$ 在 $k(A) = k$ 上的秩 (rank). 为使 $B$ 在 $A$ 上非分歧 (相应地：平展)，必要且充分条件是 $B$ 同构于 $A[t]/FA[t]$ 的商 (相应地：同构于 $A[t]/FA[t]$)，其中 $F$ 是首一可分多项式，可假设 (相应地：必然) 次数为 $n$.
\end{corollary}

\begin{proof}
只需证明必要性. 假设 $B$ 在 $A$ 上非分歧，故 $L$ 在 $k$ 上可分，由所作假设 $L/k$ 有生成元 $\xi$，于是 $\xi^i$ ($0 \leqslant i < n$) 构成 $L$ 在 $k$ 上的基. 设 $u \in B$ 提升 $\xi$，由 Nakayama 引理，$u^i$ ($0 \leqslant i < n$) 生成 $A$-模 $B$ (相应地：构成其基)，特别可以找到首一多项式 $F \in A[t]$ 使得 $F(u) = 0$，且 $B$ 同构于 $A[t]/FA[t]$ 的商 (相应地：同构于 $A[t]/FA[t]$). 最后，由应用于 $L/k$ 的 \ref{cor:I.7.4}，$F$ 和 $F'$ 模 $\mathfrak{m}_A[t]$ 生成 $A[t]$，从而 (由 $A[t]/FA[t]$ 中的 Nakayama 引理) $F$ 和 $F'$ 生成 $A[t]$，证毕.
\end{proof}

\begin{theorem}\label{thm:I.7.6}
设 $A$ 是局部环，$A \to \mathcal{O}$ 是局部同态，使得 $\mathcal{O}$ 同构于 $A$ 上有限型代数的局部化 (localized algebra). 假设 $\mathcal{O}$ 在 $A$ 上非分歧. 则可以找到 $A$ 上的整 (integral) 代数 $B$，$B$ 的一个极大理想 $n$，$B$ 在 $A$ 上的生成元 $u$，首一多项式 $F \in A[t]$，使得 $\mathfrak{r} \not\ni F'(u)$ 且 $\mathcal{O}$ (作为 $A$-代数) 同构于 $B_{\mathfrak{n}}$. 若 $\mathcal{O}$ 在 $A$ 上平展，可取 $B = A[t]/FA[t]$.
\end{theorem}

(当然，这里的条件也是充分的...)

先指出令人愉快的推论：

\begin{corollary}\label{cor:I.7.7}
为使 $\mathcal{O}$ 在 $A$ 上非分歧，必要且充分条件是 $\mathcal{O}$ 同构于一个在 $A$ 上平展的类似代数的商.
\end{corollary}

\begin{proof}
事实上，取 $\mathcal{O}' = B'_{\mathfrak{n}'}$，其中 $B' = A[t]/FA[t]$，$\mathfrak{n}'$ 是 $n$ 在 $B'$ 中的原像.
\end{proof}

\begin{corollary}\label{cor:I.7.8}
设 $f \colon X \to Y$ 是有限型态射，$x \in X$. 为使 $f$ 在 $x$ 处非分歧，必要且充分条件是存在 $x$ 的开邻域 $U$ 使得 $f|U$ 分解为 $U \to X' \to Y$，其中第一个箭头是闭浸入，第二个是平展态射.
\end{corollary}

\begin{proof}
这是 \ref{cor:I.7.7} 的简单翻译.
\end{proof}

我们来证明 \ref{thm:I.7.6} 的行话是如何从主要陈述得出的：事实上，由 \ref{cor:I.7.7} 存在满同态 $\mathcal{O}' \to \mathcal{O}$，其中 $\mathcal{O}$ 具有所要求的性质；但由于 $\mathcal{O}'$ 和 $\mathcal{O}$ 都在 $A$ 上平展，由 \ref{cor:I.4.8} 态射 $\mathcal{O}' \to \mathcal{O}$ 平展，从而是同构.

\paragraph*{\ref{thm:I.7.6} 的证明.} 它沿用 Chevalley 讨论班的一个证明. 由主定理 (Main Theorem)，有 $\mathcal{O} = B_{\mathfrak{n}}$，其中 $B$ 是 $A$ 上的有限代数，$n$ 是它的一个极大理想. 于是 $B/\mathfrak{n} = k(\mathcal{O})$ 是 $k$ 的可分扩张，因而是单生成的；若 $\mathfrak{n}_i$ ($1 \leqslant i \leqslant r$) 是 $B$ 的异于 $n$ 的极大理想，则存在 $B$ 的元素 $u$，属于所有 $\mathfrak{n}_i$，且其在 $B/\mathfrak{n}$ 中的像是该环的生成元. 而 $B/\mathfrak{n} = B_{\mathfrak{n}}/\mathfrak{n}B_{\mathfrak{n}} = B_{\mathfrak{n}}/mB_{\mathfrak{n}}$ (其中 $m$ 是 $A$ 的极大理想). 暂且承认下述引理：

\begin{lemma}\label{lem:I.7.9}
设 $A$ 是局部环，$B$ 是 $A$ 上的有限代数，$n$ 是 $B$ 的极大理想，$u$ 是 $B$ 的元素，其在 $B_{\mathfrak{n}}/mB_{\mathfrak{n}}$ 中的像作为 $k = A/\mathfrak{m}$ 上的代数生成该环，且 $u$ 属于 $B$ 的所有异于 $n$ 的极大理想. 设 $B' = B[u]$，$\mathfrak{n}' = \mathfrak{n}B'$. 则典范同态 $B'_{\mathfrak{n}'} \to B_{\mathfrak{n}}$ 是同构.
\end{lemma}

\begin{lemma}\label{lem:I.7.10}
(本应作为 \ref{prop:I.7.1} 的推论，排在它所蕴含的 \ref{cor:I.7.5} 之前.) 设 $B$ 是 $A$ 上由元素 $u$ 生成的有限代数，$n$ 是 $B$ 的极大理想，使得 $B_{\mathfrak{n}}$ 在 $A$ 上非分歧. 则存在首一多项式 $F \in A[t]$ 使得 $F(u) = 0$ 且 $F'(u) \notin \mathfrak{n}$.
\end{lemma}

\begin{proof}
事实上，设 $n$ 是 $k$-代数 $L = B \otimes_A k$ 的秩，由 Nakayama 引理存在 $A[t]$ 中的 $n$ 次首一多项式，使得 $F(u) = 0$. 设 $f$ 是由 $F$ 模 $m$ 约化得到的多项式，则 $L$ 与 $k[t]/fk[t]$ 是 $k$-同构的，故由 \ref{cor:I.7.3}，$f'(\xi)$ 不含于 $L$ 中对应于 $n$ 的极大理想 ($\xi$ 表示 $t$ 在 $L$ 中的像，即 $u$ 在 $L$ 中的像). 由于 $f'(\xi)$ 是 $F'(u)$ 的像，证毕.
\end{proof}

定理 \ref{thm:I.7.6} 现在由 \ref{lem:I.7.9} 与 \ref{lem:I.7.10} 的结合得出. 剩下要证明 \ref{lem:I.7.9}. 令 $S' = B' - \mathfrak{n}'$，故 $B'S'^{-1} = B'_{\mathfrak{n}'}$，并有交换图
\[
\begin{tikzcd}[row sep=large, column sep=large]
B \arrow[r] & BS'^{-1} \arrow[r] & BS^{-1} = B_{\mathfrak{n}} \\
B' \arrow[u] \arrow[r] & B'S'^{-1} \arrow[u] \arrow[r,"\sim"] & B'_{\mathfrak{n}'}
\end{tikzcd}
\]
同样令 $S = B - \mathfrak{n}$，故 $BS^{-1} = B_{\mathfrak{n}}$，从而有自然同态 $BS'^{-1} \to BS^{-1} = B_{\mathfrak{n}}$，我们证明它是同构，即 $S$ 的元素在 $BS'^{-1}$ 中可逆，即后者的每个极大理想 $p$ 都不与 $S$ 相交，即在 $B$ 上诱导 $n$. 事实上，由于 $BS'^{-1}$ 在 $B'S'^{-1} = B'_{\mathfrak{n}'}$ 上有限，$p$ 诱导 $B'_{\mathfrak{n}'}$ 的惟一极大理想 $\mathfrak{n}'B'_{\mathfrak{n}'}$，因而诱导 $B'$ 的极大理想 $\mathfrak{n}'$；由于 $B$ 在 $B'$ 上有限，由 $p$ 在 $B$ 上诱导的理想 $q$ 在 $\mathfrak{n}'$ 之上，必为极大理想，且不含 $u$，故等于 $n$ (这里用到了 $u$ 属于 $B$ 的所有异于 $n$ 的极大理想). 现在证明 $BS'^{-1}$ 等于 $B'S'^{-1}$：由于它在后者上有限，由 Nakayama 引理归结为证模 $\mathfrak{n}'BS'^{-1}$ 的相等，更不妨证明模 $mBS'^{-1}$ 的相等；而 $BS'^{-1}/mBS'^{-1} = B_{\mathfrak{n}}/mB_{\mathfrak{n}}$ 由 $u$ 在 $k$ 上生成 (这里用到 $u$ 的另一性质)，故 $B'$ (从而 $B'S'^{-1}$) 在其中的像是全部 (作为包含 $k$ 和 $u$ 的像的子环).

\begin{remarkstar}
定理 \ref{thm:I.7.6} 应当可以推广到 $\mathcal{O}$ 仅为半局部 (semi-local) 环的情形，以同时涵盖 \ref{cor:I.7.5}：我们假设 $\theta/m\mathcal{O}$ 是单生成的 $k$-代数；于是可以找到 $u \in B$，其在 $B/mB$ 中的像是生成元，且属于 $B$ 的所有不由 $\mathcal{O}$ 诱导的极大理想. 引理 \ref{lem:I.7.9} 和 \ref{lem:I.7.10} 应当不难适应. 更一般地，...
\end{remarkstar}

\section{平展概形的无穷小提升。形式概形的应用}

\begin{proposition}\label{prop:I.8.1}
设 $Y$ 是预概形，$Y_0$ 是子预概形，$X_0$ 是平展 $Y_0$-概形，$x$ 是 $X_0$ 的点. 则存在平展 $Y$-概形 $X$，$x$ 在 $X_0$ 中的邻域 $U_0$，以及 $Y_0$-同构
\[
U_0 \stackrel{\sim}{\longrightarrow} X \times_Y Y_0.
\]
\end{proposition}

\begin{proof}
事实上，设 $y$ 是 $x$ 在 $Y_0$ 中的投影，把 \ref{thm:I.7.6} 应用于 $Y_0$ 和 $X_0$ 中点 $y$ 与 $x$ 的局部环的平展局部同态 $A_0 \to B_0$：得到同构
\[
B_0 = C_{0,\mathfrak{n}_0}, \qquad C_0 = A_0[t]/F_0 A_0[t],
\]
其中 $F_0$ 是首一多项式，$\mathfrak{n}_0$ 是 $C_0$ 的一个不含 $F'_0(t)$ 在 $C_0$ 中的类的极大理想. 设 $A$ 是 $y$ 在 $Y$ 中的局部环，$F$ 是 $A[t]$ 中的首一多项式，经满同态 $A \to A_0$ 给出 $F_0$ (提升 $F_0$ 的系数)，最后设 $C = A[t]/FA[t]$，$n$ 是 $C$ 中由自然满同态 $C \to C \otimes_A A_0 = C_0$ 下 $n$ 的原像构成的极大理想. 令
\[
B = C_{\mathfrak{n}}.
\]
由构造和 \ref{prop:I.7.1} 立即可知 $B$ 在 $A$ 上平展，且有同构 $B \otimes_A A_0 = A_0$ (原文如此，应为 $B_0$). 由第一章 (Chap. I) 知道，存在有限型 $Y$-概形 $X$ 和 $X$ 中位于 $y$ 之上的点 $z$，使得 $\mathcal{O}_z$ 与 $C$ 是 $A$-同构的；由于后者在 $A = \mathcal{O}_y$ 上平展，可以 (取 $X$ 足够小) 假设 $X$ 在 $Y$ 上平展. 设 $X'_0 = X \times_Y Y_0$，则 $z$ 在 $X'_0$ 中的局部环等同于 $\theta_z \otimes_A A_0 = B \otimes_A A_0$，故同构于 $B_0$. 这个同构由一个把 $x$ 在 $X_0$ 中的邻域 $U_0$ 映到 $z$ 在 $X'_0$ 中的邻域上的同构 (上引文献 (loc. cit.)) 定义，取 $X$ 足够小可以假设该邻域等于 $X'_0$. 证毕.
\end{proof}

\begin{corollary}\label{cor:I.8.2}
对平展覆盖有类似的陈述，需假设剩余域 $k(y)$ 无限.
\end{corollary}

\begin{proof}
证明相同，以 \ref{cor:I.7.5} 代替 \ref{thm:I.7.6}.
\end{proof}

\begin{theorem}\label{thm:I.8.3}
\ref{thm:I.5.5} 中所考虑的函子是范畴的等价.
\end{theorem}

\begin{proof}
由定理 \ref{thm:I.5.5}，剩下要证明每个平展 $S_0$-概形 $X_0$ 都同构于某个 $X \times_S S_0$，其中 $X$ 是平展 $S$-概形. $X$ 的底空间 (underlying topological space) 必然与 $X_0$ 的相同，且 $X_0$ 还等同于 $X$ 的闭子预概形. 因此问题等价于：在 $X_0$ 的底空间 $|X_0|$ 上找 $f_0^{*}(\mathcal{O}_S)$ 的代数层 $\mathcal{O}_X$ (其中 $f_0$ 是投影 $X_0 \to S_0$，这里看作底空间之间的连续映射)，使 $|X_0|$ 成为平展 $S$-预概形 $X$，并有代数同态 $\mathcal{O}_X \to \mathcal{O}_{X_0}$，与标量层上的同态 $f_0^{*}(\mathcal{O}_S) \to f_0^{*}(\mathcal{O}_{S_0})$ 相容，诱导同构 $\mathcal{O}_X \otimes_{f_0^{*}(\mathcal{O}_S)} f_0^{*}(\mathcal{O}_{S_0}) \stackrel{\sim}{\longrightarrow} \mathcal{O}_{X_0}$ (于是 $X$ 将是沿 $X_0$ 约化的平展 $S$-预概形，因 $X_0$ 在 $S_0$ 上分离故 $X$ 在 $S$ 上分离，且 $X$ 回答了问题). 此外，若 $(U_i)$ 是 $X_0$ 的开覆盖，且在每个 $U_i$ 上都找到了问题的解，则由唯一性定理 \ref{thm:I.5.5}，这些解可以粘合 (i.e. 定义它们的代数层连同它们的增广同态可以粘合)，并且立即看到这样在 $S$ 之上构造的环空间 (ringed space) 是带同构 $X \times_S S_0 \to X_0$ 的平展 $S$-预概形 $X$. 因此只需局部求解，这由 \ref{prop:I.8.1} 保证.
\end{proof}

\begin{corollary}\label{cor:I.8.4}
设 $S$ 是局部诺特形式预概形，带定义理想 (ideal of definition) $J$，设 $S_0 = (|S|, \mathcal{O}_S/\mathcal{J})$ 是相应的普通预概形. 则函子 $\mathfrak{X} \mapsto \mathfrak{X} \times_S S_0$ 从 $S$ 的平展覆盖的范畴到 $S_0$ 的平展覆盖的范畴是范畴的等价.
\end{corollary}

当然，形式预概形 $S$ 的平展覆盖是指 $S$ 的覆盖，i.e. $S$ 上的形式预概形，由凝聚代数层 $\mathcal{B}$ 定义，使得 $\mathcal{B}$ 局部自由，且 $\mathcal{B}$ 的剩余纤维 $\mathcal{B}_s \otimes_{\mathcal{O}_s} k(s)$ 是 $k(s)$ 上的可分代数. 若以 $S_n$ 记普通预概形 $\bigl(|S|, \mathcal{O}_S/J^{n+1}\bigr)$，则 $S$ 上凝聚代数层 $\mathcal{B}$ 的给出等价于诸 $S_n$ 上一族凝聚代数层 $\mathcal{B}_n$ 的给出，它们带有传递的同态系统 $\mathcal{B}_m \to \mathcal{B}_n$ ($m \geqslant n$)，定义同构 $\mathcal{B}_m \otimes_{\mathcal{O}_{S_m}} \mathcal{O}_{S_n} \stackrel{\sim}{\longrightarrow} \mathcal{B}_n$. 立即可知 $\mathcal{B}$ 局部自由当且仅当诸 $S_n$ 上的 $\mathcal{B}_n$ 局部自由，且可分性条件成立当且仅当它对 $\mathcal{B}_0$ 成立，或对所有的 $\mathcal{B}_n$ 成立. 这样，$\mathcal{B}$ 在 $S$ 上平展当且仅当诸 $S_n$ 上的 $\mathcal{B}_n$ 平展. 据此，\ref{cor:I.8.4} 立即由 \ref{thm:I.8.3} 得出.

\begin{remarkstar}
在 \ref{cor:I.8.4} 中没有必要局限于覆盖的情形. 不过目前只用到这一情形.
\end{remarkstar}

\section{永久性质}

设 $A \to B$ 是局部且平展的同态，这里我们考察若干``对 $A$ 成立的某种性质蕴含对 $B$ 的同一性质成立，或反之''的情形. 其中相当一部分命题只是 $B$ 在 $A$ 上拟有限且平坦这一简单事实的推论，我们只限于``回忆''其中的一些：$A$ 与 $B$ 有相同的 Krull 维数 (Krull dimension)，相同的深度 (depth) (在仍通用的术语中即 Serre 的``上同调余维数'' (cohomological codimension)). 例如由此得出 $A$ 是 Cohen-Macaulay 的当且仅当 $B$ 是. 此外，对 $B$ 的每个在 $A$ 上诱导 $p$ 的素理想 $q$，$B_{\mathfrak{q}}$ 仍在 $A_{\mathfrak{p}}$ 上拟有限且平坦，只要假设 $B$ 是 $A$ 上有限型代数的局部化 (这由``有限型态射为拟有限 (相应地：平坦) 的点集是开的''这一事实得出)；而且 $A$ 的每个素理想 $p$ 都由 $B$ 的某个素理想 $q$ 诱导 (因为 $B$ 在 $A$ 上忠实平坦). 例如由此得出 $p$ 与 $q$ 有相同的秩 (rank)；还得出 $A$ 无嵌入素理想 (embedded prime ideal) 当且仅当 $B$ 无嵌入素理想.

因此我们只限于对平展态射更特殊的情形.

\begin{proposition}\label{prop:I.9.1}
设 $A \to B$ 是局部平展同态. 为使 $A$ 正则 (regular)，必要且充分条件是 $B$ 正则.
\end{proposition}

\begin{proof}
事实上，设 $k$ 是 $A$ 的剩余域，$L$ 是 $B$ 的剩余域. 由于 $B$ 在 $A$ 上平坦且 $L = B \otimes_A k$，i.e. $n = mB$ (其中 $m$，$n$ 是 $A$，$B$ 的极大理想)，$B$ 上的 $m$-adic 滤链 (filtration) 与其 $n$-adic 滤链相同，且有
\[
\operatorname{gr}^{*}(B) = \operatorname{gr}^{*}(A) \otimes_k L.
\]
由此推出 $\mathrm{gr}^{*}(B)$ 是 $L$ 上的多项式代数当且仅当 $\mathrm{gr}^{*}(A)$ 是 $k$ 上的多项式代数. 证毕. (N.B. 没有用到 $L/k$ 可分这一事实).
\end{proof}

\begin{corollary}\label{cor:I.9.2}
设 $f \colon X \to Y$ 是平展态射. 若 $Y$ 正则\index[term]{pingzhantaishe@平展态射}，则 $X$ 正则，逆命题在 $f$ 满射时成立.
\end{corollary}

% 原书编号笔误: 下面这一条在原文中与上面的推论同为 ``\ref{cor:I.9.2}''
\begingroup
\addtocounter{theorem}{-1}%
\renewcommand{\thetheorem}{\thesection.2}%
\begin{proposition}\label{prop:I.9.2red}
设 $f \colon X \to Y$ 是平展态射. 若 $Y$ \emph{约化} (reduced)，则 $X$ 亦然，逆命题在 $f$ 满射时成立.
\end{proposition}

\endgroup

这等价于：

\begin{corollary}\label{cor:I.9.3}
设 $f \colon A \to B$ 是局部平展同态，$B$ 同构于 $A$ 上有限型代数的局部化. 为使 $A$ 约化，必要且充分条件是 $B$ 约化.
\end{corollary}

\begin{proof}
必要性是平凡的，因为 $A \to B$ 单射 ($B$ 在 $A$ 上忠实平坦). 充分性：设 $\mathfrak{p}_i$ 是 $A$ 的极小素理想，由假设自然映射 $A \to \prod_i A/\mathfrak{p}_i$ 单射，用平坦 $A$-模 $B$ 作张量，得到 $B \to \prod_i B/\mathfrak{p}_i B$ 单射，归结为证明诸 $B/\mathfrak{p}_i B$ 约化. 由于 $B/\mathfrak{p}_i B$ 在 $A/\mathfrak{p}_i$ 上平展，归结到 $A$ 整 (integral domain) 的情形. 设 $K$ 是其分式域 (field of fractions)，则 $A \to K$ 单射，故 ($B$ 为 $A$-平坦) $B \to B \otimes_A K$ 亦单射，归结为证明后一个环约化. 而 $B$ 是 $A$ 上有限型代数的局部化，它是 $Y = \mathrm{Spec}(A)$ 上有限型平展概形 $X = \mathrm{Spec}(C)$ 中某点 $x$ 的局部环，故 $B \otimes_A K$ 是 $X \otimes_A K$ 的环 $C \otimes_A K$ 的局部化 (关于某个适当的乘法稳定集 (multiplicatively closed set)). 由于 $X \otimes_A K$ 在 $K$ 上平展，其环是有限个域的乘积 ($K$ 的可分扩张)，$B \otimes_A K$ 亦然，证毕.
\end{proof}

\begin{corollary}\label{cor:I.9.4}
设 $f \colon A \to B$ 是局部平展同态，假设 $A$ 解析约化 (analytically reduced) (i.e. $A$ 的完备化 $\widehat{A}$ 无幂零元 (nilpotent element)). 则 $B$ 解析约化，从而是约化的.
\end{corollary}

\begin{proof}
事实上，$\widehat{B}$ 在 $\widehat{A}$ 上有限且平展，应用 \ref{cor:I.9.3} 即得.
\end{proof}

\begin{theorem}\label{thm:I.9.5}
设 $f \colon A \to B$ 是局部同态，$B$ 同构于 $A$ 上有限型代数的局部化. 则：
\begin{enumerate}[(i)]
\item 若 $f$ 平展，则 $A$ 正规 (normal) 当且仅当 $B$ 正规.
\item 若 $A$ 正规，则 $f$ 平展当且仅当 $f$ 单射且非分歧 (此时由 (i)，$B$ 正规).
\end{enumerate}
\end{theorem}

\begin{proof}
我们对 (i) 给出两个不同的证明，第一个用拟有限平坦态射的某些性质 (本节开头所回忆的)，不用 \ref{thm:I.7.6} (从而不用主定理)；第二个证明则相反. 最后，对 (ii) 看来无论如何都需要主定理.
\end{proof}

\paragraph*{第一个证明.} 使用下述诺特局部环 $A$ (维数 $\neq 0$) 正规的必要充分条件.

\begin{remarkstar}
\textbf{\emph{Serre 判则} (Serre criterion).}\index[term]{Serre panze@Serre 判则} (i) 对 $A$ 的每个秩 1 的素理想 $p$，$A_{\mathfrak{p}}$ 正规 (等价地：正则)；(ii) 对 $A$ 的每个秩 $\geqslant 2$ 的素理想 $p$，深度 $A_{\mathfrak{p}} \geqslant 2$\footnote{参见 EGA IV 5.8.6.}.
\end{remarkstar}

这里我们承认这个判则，它本应出现在关于平坦态射的那一节. 它的主要优点是它不预先假设 $A$ 约化，更不假设 $A$ 整. 这里已经可以假设 $\dim A = \dim B \neq 0$.

由本节开头的回忆，$A$ 中秩 1 (相应地：秩 $\geqslant 2$) 的素理想 $p$ 恰是 $B$ 中秩 1 (相应地：秩 $\geqslant 2$) 的素理想 $q$ 在 $A$ 上的迹. 最后，若 $p$ 与 $q$ 对应，则 $B_{\mathfrak{q}}$ 在 $A_{\mathfrak{p}}$ 上平展，故与 $A_{\mathfrak{p}}$ 有相同的深度，且正则当且仅当 $A_{\mathfrak{p}}$ 正则 (\cref{prop:I.9.1}). 应用 Serre 判则，得到 $A$ 正规当且仅当 $B$ 正规.

\paragraph*{第二个证明.} 假设 $B$ 正规，设 $L$ 是其分式域，$K$ 是 $A$ 的分式域 ($A$ 因 $B$ 整而整). 在 \ref{cor:I.9.3} 的证明中看到 $B \otimes_A K$ 是有限个域的直积，由于它含于 $L$ 故是域，又因它包含 $B$ 故等于 $L$. $K$ 中在 $A$ 上整的元素在 $B$ 上整，因 $B$ 正规故属于 $B$，又因 $B \cap K = A$ (由 $B$ 在 $A$ 上忠实平坦得出) 故属于 $A$.

现在假设 $A$ 正规，证明 $B$ 正规. 由 \ref{thm:I.7.6} 有 $B = B'_{\mathfrak{n}}$，其中 $B' = A[t]/FA[t]$，$F$ 和 $n$ 如 \ref{thm:I.7.6} 所述. 于是 $\mathcal{L} = \mathcal{B} \otimes_A K$ 是 $B' \otimes_A K = K[t]/FK[t]$ 的局部化，并且是域的直积；这些 $K$ 的有限可分扩张 (每当我们把 Artin 环，这里是 $B'_K$，关于一个乘法稳定集局部化时) 是 $B'_K$ 的直因子，因而对应于 $K[t]$ 中的分解 $F = F_1 F_2$，而 $L$ 中对应于 $t$ 的生成元已被 $F_1$ 零化. 再者，因 $A$ 正规，$F_i \in A[t]$ (假设它们首一). 注意 $B \to L = B \otimes_A K$ 单射 ($A \to K$ 单射且 $B$ 在 $A$ 上平坦)，故已有 $F_1(u) = 0$，其中 $u$ 是 $t$ 在 $L$ 中的类. 假设 $F$ 的次数取最小，则 $F_2 = 1$ (N.B. 有 $F'(u) = F'_1(u)F_2(u) + F_1(u)F'_2(u) = F'_1(u)F_2(u)$ 因为 $F_1(u) = 0$，故由 $F'(u) \neq 0$ 得 $F'_1(u) \neq 0$.)

因此有
\begin{equation}
L = B \otimes_A K = K[t]/FK[t] \tag{($*$)}
\end{equation}
$F$ 因而是 $K[t]$ 中的可分多项式 (但显然不必在 $A[t]$ 中). (N.B. 目前我们本质上只证明了：在 \ref{thm:I.7.6} 中可以选取 $F$ 和 $n$，使得 —— 用这里的记号 —— $B' \to B'_{\mathfrak{n}} = B$ 单射；这用到了 $A$ 的正规性；我不知道无正规性假设时是否仍然成立).

现在回忆一个众所周知的引理，取自 Serre 去年的课程：

\begin{lemma}\label{lem:I.9.6}
设 $K$ 是环，$F \in K[t]$ 是首一可分多项式，$L = K[t]/FK[t]$，$u$ 是 $t$ 在 $L$ 中的类 (从而 $F'(u)$ 是 $L$ 的可逆元). 则有以下公式 (其中 $n = \deg F$)：
\[
\mathrm{tr}_{L/K}\, u^{i}/F'(u) = 0 \quad \text{若 } 0 \leqslant i < n - 1,
\]
\[
\mathrm{tr}_{L/K}\, u^{n-1}/F'(u) = 1.
\]
\end{lemma}

\begin{corollary}\label{cor:I.9.7}
矩阵 $\bigl(u^{j} \cdot u^{i}/F'(u)\bigr)_{0 \leqslant i, j \leqslant n-1}$ 的行列式等于 $(-1)^{n(n-1)/2}$，因而在 $K$ 的任何子环 $A$ 中可逆.
\end{corollary}

\begin{corollary}\label{cor:I.9.8}
设 $A$ 是 $K$ 的子环，$V$ 是 $L$ 中由 $u^i$ ($0 \leqslant i \leqslant n-1$) 生成的 $A$-模，$V'$ 是 $L$ 中由满足``对一切 $y \in V$ (i.e. 对形如 $u^i$，$0 \leqslant i \leqslant n-1$ 的 $y$) 有 $\mathrm{tr}_{L/K}(xy) \in A$''的 $x \in L$ 构成的子 $A$-模. 则 $V'$ 是以 $u^{i}/F'(u)$ ($0 \leqslant i \leqslant n-1$) 为基的 $A$-模.
\end{corollary}

\begin{corollary}\label{cor:I.9.9}
假设 $K$ 是整正规环 $A$ 的分式域，$F$ 的系数在 $A$ 中. 则用 \ref{cor:I.9.8} 的记号，$V'$ 包含 $A$ 在 $L$ 中的正规闭包 (normal closure) $A'$，后者因而包含在 $A[u]/F'(u)$ 中，从而更包含在 $A[u][F'(u)^{-1}]$ 中.
\end{corollary}

\begin{proof}
把这个最后的推论应用于我们证明中得到的情形：由于 $F'(u)$ 在包含 $A[u]$ 的 $B$ 中可逆，$B$ 包含 $A'$. 由主定理 (或由 $B = A[u]_{\mathfrak{n}}$ 这一事实)，$B$ 是 $A'$ 的局部化代数. 由于 $A'$ 正规，$B$ 亦然.
\end{proof}

\paragraph*{(ii) 的证明.} 像前面的证明那样，证明在 \ref{thm:I.7.6} 中可以选取 $F$ 使得 ($*$) 仍成立. 惟一的先行障碍是，不再假设 $B$ 在 $A$ 上平坦，就不能再断言 $B \to L$ 单射，因此推理原则上只适用于 $B$ 在该同态下的像 $B_1$. 由此立即得出 $B_1$ 在 $A$ 上平坦 (作为 $A$ 上自由代数的局部化). 由 \ref{cor:I.4.8}，态射 $B \to B_1$ 平展，从而是同构，证明完毕.

(从写作的角度看，应当把最后两个证明的顺序对调，并把引理及其推论的形式计算放到单独的一节.)

\begin{corollary}\label{cor:I.9.10}
设 $f \colon X \to Y$ 是平展态射. 若 $Y$ 正规，则 $X$ 正规，逆命题在 $f$ 满射时成立.
\end{corollary}

\begin{corollary}\label{cor:I.9.11}
设 $f \colon X \to Y$ 是支配 (dominant) 态射，$Y$ 正规且 $X$ 连通. 若 $f$ 非分歧，则 $f$ 平展，故 $X$ 正规，从而 (连通) 不可约 (irreducible).
\end{corollary}

\begin{proof}
设 $U$ 是 $f$ 平展的点集，它是开的，只需证明它也是闭的且非空. $U$ 包含 $Y$ 的生成点 (generic point) 的原像 (因为对域上的代数，非分歧 = 平展)，故 ($X$ 支配 $Y$) 非空. 若 $x$ 属于 $U$ 的闭包，则它属于 $U$ 的某个不可约分支 $U_i$ 的闭包，因而属于 $X$ 的与 $U_i$ 相交的某个不可约分支 $X_i = \overline{U}_i$，从而支配 $Y$ (因为 $U$ 的每个分支都在 $Y$ 上平坦，因而支配 $Y$). 于是，若 $y$ 是 $x$ 在 $Y$ 上的投影，则 $\mathcal{O}_y \to \mathcal{O}_x$ 单射 (考虑到 $\mathcal{O}_y$ 整). 由于 $\mathcal{O}_y$ 正规且 $\mathcal{O}_y \to \mathcal{O}_x$ 非分歧，借助 \ref{thm:I.9.5}(ii) 即得结论.
\end{proof}

\begin{corollary}\label{cor:I.9.12}
设 $f \colon X \to Y$ 是有限型支配态射，$Y$ 正规且 $X$ 不可约. 则 $f$ 平展的点集恰为 $\varOmega_{X/Y}^{1}$ 的支集 (support) 的补集，i.e. 由差理想 $\mathfrak{d}_{X/Y}$ 定义的 $X$ 的闭子预概形的补集.
\end{corollary}

(这正是 §4 的注记中所提到的``不那么平凡的''陈述.)

\begin{remarkstar}
我们切勿以为不可约概形的连通平展覆盖本身不可约，当不假设基正规时尤其如此. 这个问题将在 §11 中研究.
\end{remarkstar}

\section{正规概形的平展覆盖}

\begin{proposition}\label{prop:I.10.1}
设 $X$ 是正规连通概形 $Y$ (其函数域为 $K$) 上的平展分离预概形. 则 $X$ 的连通分支 $X_i$ 是整的，其函数域 $K_i$ 是 $K$ 的有限可分扩张，且 $X_i$ 等同于 $X$ 在 $K_i$ 中的正规化 (normalized) 的一个非空开部分 (故 $X$ 等同于 $Y$ 在 $R(X) = L = \prod K_i$ 中的正规化的一个开稠密部分).
\end{proposition}

\begin{proof}
由 \ref{cor:I.9.10}，$X$ 正规，更不用说其局部环整，故 $X$ 的连通分支不可约. 由于 $X_i$ 正规，且在 $Y$ 上有限且支配，由主定理的一个 (事实上几乎平凡) 特殊情形得出，$X_i$ 是 $X$ 在 $K_i$ 中的正规化的一个开集.
\end{proof}

\begin{corollary}\label{cor:I.10.2}
在 \ref{prop:I.10.1} 的条件下，$X$ 在 $Y$ 上有限 (i.e. 是 $Y$ 的平展覆盖) 当且仅当 $X$ 同构于 $Y$ 在 $L = R(X)$ (即 $X$ 上的有理函数 (rational functions) 环) 中的正规化 $Y'$.
\end{corollary}

\begin{proof}
事实上，我们知道这个正规化在 $Y$ 上有限 ($Y$ 正规且 $R/K$ 可分)，反之若 $X$ 在 $Y$ 上有限则它在 $Y'$ 上有限，故其在 $Y'$ 中的像是闭的，另一方面它又是稠密的.
\end{proof}

$K$ 上的有限阶代数 $L$ 称为在 $X$ 上非分歧 (或简称在 $K$ 上非分歧，若 $X$ 隐含) 的，若 $L$ 是 $K$ 上的可分代数，i.e. 是有限个可分扩张 $K_i$ 的直积，且 $Y$ 在 $L$ 中的正规化 $Y'$ ($Y$ 在诸 $K_i$ 中的正规化的不交并) 在 $Y$ 上非分歧 ($=$ 由 \ref{cor:I.9.11} 平展). 因此：

\begin{corollary}\label{cor:I.10.3}
对每个在 $Y$ 上有限、其每个不可约分支都支配 $Y$ 的 $X$，设 $R(X)$ 是 $X$ 上的有理函数环 ($X$ 的不可约分支的生成点的局部环的乘积)，于是 $X \mapsto R(X)$ 是取值于 $K = R(Y)$ 上有限阶代数的函子. 该函子建立 $Y$ 的连通平展覆盖的范畴与 $K$ 的在 $Y$ 上非分歧的扩张 $L$ 的范畴之间的等价.
\end{corollary}

逆函子是正规化函子.

假设 $Y$ 仿射，故由分式域为 $K$ 的正规环 $A$ 定义\index[term]{zhengguigaixing@正规概形}. 设 $L$ 是 $K$ 的有限扩张，是域的直积，则按定义 $Y$ 在 $L$ 中的正规化 $Y'$ 同构于 $\mathrm{Spec}(A')$，其中 $A'$ 是 $A$ 在 $L$ 中的正规化. 说 $L$ 在 $Y$ 上非分歧，意味着 $A'$ 在 $A$ 上非分歧 (或：平展). 若 $A$ 局部，则等价于说局部环 $A'_{\mathfrak{n}}$ ($n$ 遍历 $A'$ 的有限个极大理想，即诱导 $A$ 的极大理想 $m$ 的素理想) 在局部环 $A$ 上非分歧 ($=$ 平展). 最后还注意到，判别式判则 (\cref{prop:I.4.10}) 在这一情形也可应用 (更一般地，该判则的一个变体应当这样陈述，无需预先的平坦性条件，当 $X$ 支配 $Y$，$Y$ 仍假设局部整时：$A \to B$ 和 $B \to B \otimes_A K = L$ 单射 —— 则 $\mathrm{tr}_{L/K}$ 有定义 —— 且 $\mathrm{tr}_{L/K}(xy)$ 诱导基本双线性型 $B \times B \to A$，i.e. 存在 $x_i \in B$ ($1 \leqslant i \leqslant n$，$n$ 为 $L$ 在 $K$ 上的秩) 使得对所有 $i, j$ 有 $\mathrm{tr}(x_i x_j) \in A$ 且 $\det\bigl(\mathrm{tr}(x_i x_j)\bigr)$ 在 $A$ 中可逆).

推论组 (\cref{prop:I.4.6}) 立即蕴含经典框架中非分歧性的推论组：

\begin{proposition}\label{prop:I.10.4}
设 $Y$ 是正规整预概形，函数域为 $K$. (i) $K$ 在 $Y$ 上非分歧. (ii) 若 $L$ 是 $K$ 的在 $Y$ 上非分歧的扩张，$Y'$ 是函数域为 $L$ 且支配 $Y$ 的正规预概形 (例如 $Y$ 在 $L$ 中的正规化)，$M$ 是 $L$ 的在 $Y'$ 上非分歧的扩张，则 $M/K$ 在 $X$ 上非分歧 (非分歧性的传递性 (transitivity)). (iii) 设 $Y'$ 是支配 $Y$ 的正规整预概形，函数域为 $K'/K$；若 $L$ 是 $K$ 的在 $Y$ 上非分歧的扩张，则 $L \otimes_K K'$ 是 $K'$ 的在 $Y'$ 上非分歧的扩张 (平移性质 (translation property)).
\end{proposition}

此外：

\begin{corollary}\label{cor:I.10.5}
在 (iii) 的条件下，若 $Y = \mathrm{Spec}(A)$，$Y' = \mathrm{Spec}(A')$，则 $Y'$ 在 $L' = L \otimes_K K'$ 中的正规化 $\overline{A'}$ 等同于 $\overline{A} \otimes_A A'$，其中 $\overline{A}$ 是 $A$ 在 $L$ 中的正规化.
\end{corollary}

通常人们 (那些厌恶考虑非整环、即使它是域的直积的人) 把平移性质表述为如下 (更弱) 的形式：

\begin{corollary}\label{cor:I.10.6}
在 (iii) 的条件下，设 $L_1$ 是 $L/K$ (在 $Y$ 上非分歧) 与 $K'/K$ 的复合扩张 (composite extension). 则 $L_1/K'$ 在 $Y'$ 上非分歧. 在 $Y = \mathrm{Spec}(A)$，$Y' = \mathrm{Spec}(A')$ 的情形，还成立
\[
\overline{A'} = A[\overline{A}, A']
\]
i.e. $A'$ 在 $L_1$ 中的正规化环 $\overline{A'}$ 是由 $A'$ 和 $A$ 在 $L$ 中的正规化 $\overline{A}$ 生成的 $A$-代数.
\end{corollary}

后一事实在无非分歧性假设时是假的，即使在数域 (number field) 的复合扩张的情形也是如此...

为了结束本节，我们给出与平展覆盖概念的直观图像相对应的解释：在 $Y$ 的所考虑点 $y$ 之上应当有``最大数目''的点，特别地，在 $y$ 之上不应当有``多个重合的点''. 为了以所希望的全部普遍性证明这个方向的结果，我们在这里承认下面的命题 \ref{prop:I.10.7} (其证明将在大全第四章 §15 中给出，它使用 Chevalley 的可构造集 (constructible sets) 技术\index[term]{kegoubufen@可构部分}，以及一点点下降 (descent) 理论...).

有限型态射 $f \colon X \to Y$ 称为泛开 (universally open)\index[term]{wanyoukaitaishe@万有开态射} 的，若对任何基扩张 $Y' \to Y$ (其中 $Y'$ 局部诺特)，态射 $f' \colon X' = X \times_Y Y' \to Y'$ 是开的，i.e. 把开集变成开集. 还可以只限于 $Y'$ 在 $Y$ 上有限型的情形 (甚至 $Y'$ 取形如 $Y[t_1, \dots, t_r]$，其中 $t_i$ 是未定元). 泛开态射更是开的 (逆命题不真)，另一方面若 $f$ 开，且 $X$ 与 $Y$ 不可约，则 $f$ 的所有纤维的所有分支都有相同的维数 (即生成纤维 $f^{-1}(z)$ 的维数，$z$ 为 $Y$ 的生成点). 最后，若 $Y$ 正规，则最后一个条件已经蕴含 $f$ 泛开 (Chevalley 定理). 例如由此得出：若 $f \colon X \to Y$ 是拟有限态射，$Y$ 正规不可约，则 $f$ 泛开 (或：开) 当且仅当 $X$ 的每个不可约分支都支配 $Y$. 还回忆：(有限型的) 平坦态射既然开，也就泛开. 在这些准备之后，``回忆''下述

\begin{proposition}\label{prop:I.10.7}
设 $f \colon X \to Y$ 是拟有限分离泛开态射. 对每个 $y \in Y$，设 $n(y)$ 是纤维 $f^{-1}(y)$ 的``点的几何个数''\index[term]{jihedianshu@几何点数} (geometric number of points)，等于对 $x \in f^{-1}(y)$，剩余扩张 $k(x)/k(y)$ 的可分次数 (separable degree) 之和. 则 $Y$ 上的函数 $y \mapsto n(y)$ 上半连续 (upper semicontinuous). 为使它在点 $y$ 的邻域上为常数 (i.e. 使 $n(y) = n(z_i)$，其中 $z_i$ 是 $Y$ 的包含 $y$ 的不可约分支的生成点)，必须存在 $y$ 的邻域 $U$ 使得 $X|U$ 在 $U$ 上有限\footnote{参见 EGA IV 13.1.4 (原书注 (8)).}.
\end{proposition}

\begin{corollary}\label{cor:I.10.8}
若 $y \mapsto n(y)$ 为常数且 $Y$ \emph{几何单支} (geometrically unibranched)\index[term]{jihedanzhi@几何单支}\footnote{参见 §11 的末尾 (原书注 (9)).}，则 $X$ 的不可约分支两两不交.
\end{corollary}

\begin{proposition}\label{prop:I.10.9}
设 $f \colon X \to Y$ 是平展分离态射. 用 \ref{prop:I.10.7} 的记号，函数 $y \mapsto n(y)$ 上半连续. 为使它在点 $y$ 的邻域上为常数 (i.e. 使 $n(y) = n(z_i)$，其中 $z_i$ 是 $Y$ 的包含 $y$ 的不可约分支的生成点)，必要且充分条件是存在 $y$ 的开邻域 $U$ 使得 $X|U$ 在 $U$ 上有限，i.e. 是 $U$ 的平展覆盖.
\end{proposition}

\begin{corollary}\label{cor:I.10.10}
为使平展分离态射 $f \colon X \to Y$ ($Y$ 连通) 有限 (i.e. 使 $X$ 成为 $Y$ 的平展覆盖)，必要且充分条件是 $f$ 的所有纤维有相同的点的几何个数.
\end{corollary}

在 \ref{prop:I.10.7} 及其推论中，对 $Y$ 没有正规性假设. 若作这样的假设，就得到更强的陈述 (通常被取作覆盖非分歧性的定义)：

\begin{theorem}\label{thm:I.10.11}
设 $f \colon X \to Y$ 是拟有限分离态射. 假设 $Y$ 不可约，$X$ 的每个分支都支配 $Y$，$X$ 约化 (i.e. $\mathcal{O}_X$ 无幂零元). 设 $n$ 是 $X$ 在 $Y$ 上的次数 (degree) ($X$ 的不可约分支 $X_i$ 的域 $K_i$ 在 $Y$ 的域 $K$ 上的次数之和). 设 $y$ 是 $Y$ 的正规点. 则 $X$ 在 $y$ 之上的点的几何个数 $n(y)$ 小于等于 $n$，等号成立当且仅当存在 $y$ 的开邻域 $U$ 使得 $X|U$ 是 $U$ 的平展覆盖.
\end{theorem}

\begin{proof}
``仅当''是平凡的，证明``当''. 设 $z$ 是 $Y$ 的生成点，有 $n(z) =$ (诸 $K_i/K$ 的可分次数之和) $\leqslant n$，由 \ref{prop:I.10.7} 有 $n(y) \leqslant n(z) \leqslant n$，等号蕴含 $X|U$ 对 $y$ 的某个适当邻域 $U$ 在 $U$ 上有限. 故可假设 $X$ 在 $Y$ 上有限，且 $Y$ 上的函数 $n(y')$ 为常数. 最后由 \ref{cor:I.10.8}，$X$ 是其不可约分支的不交并，为证明它在 $y$ 处非分歧，归结到 $X$ 不可约从而整的情形. 最后可以假设 $Y = \mathrm{Spec}(\mathcal{O}_y)$. 于是定理归结为下述经典陈述：
\end{proof}

\begin{corollary}\label{cor:I.10.12}
设 $A$ 是 (如往常一样诺特的) 正规局部环，函数域为 $K$，$L$ 是 $K$ 的 $n$ 次有限扩张，可分次数为 $n_s$，$B$ 是 $L$ 的包含在 $L$ 中的子环，在 $A$ 上有限，分式域为 $L$，$m$ 是 $A$ 的极大理想，$n'$ 是 $B/\mathfrak{m}B$ 在 $A/\mathfrak{m}A = k$ 上的可分次数 (= 该环的剩余扩张的可分次数之和). 有 $n' \leqslant n_s$，从而 $n' \leqslant n$. 最后一个不等式取等号当且仅当 $B$ 在 $A$ 上非分歧 ($=$ 平展).
\end{corollary}

\begin{proof}
剩下只需证明 $n' = n$ 蕴含 $B$ 在 $A$ 上平展. 回忆 $k$ 无限时的证明：只需证明 $R = B/\mathfrak{m}B$ 在 $k$ 上可分；若不然，则 (由已知引理) 存在 $R$ 的元素 $a$，其在 $k$ 上的最小多项式次数 $> n'$. 该元素来自 $B$ 的元素 $x$，后者在 $K$ 上 (作为 $L$ 的元素) 的最小多项式次数 $\leqslant n$；另一方面，由于 $A$ 正规，该最小多项式的系数在 $A$ 中，模 $m$ 约化给出 $k[t]$ 中的首一多项式 $F$，次数 $\leqslant n = n'$，使得 $F(a) = 0$，矛盾.
\end{proof}

在一般情形 ($k$ 可以有限)，回到几何语言，考虑 $Y' = \mathrm{Spec}(A[t])$，它在 $Y$ 上忠实平坦，以及 $Y'$ 在 $y$ 上的纤维 $\mathrm{Spec}(k[t])$ 的生成点 $y'$. 则 $X$ 在 $Y$ 上的 $y$ 处非分歧，当且仅当 $X' = X \times_Y Y' = \mathrm{Spec}(B[t])$ 在 $Y'$ 上的 $y'$ 处非分歧，这立即可见. 另一方面，由 $y'$ 的选取，其剩余域是 $k(t)$，因而是无限的. 由于 $y'$ 是 $Y'$ 的正规点，归结到前面的情形.

\section{一些补充}

我们已经说过，整概形的连通平展覆盖未必整. 下面是这一事实的两个例子.

\begin{enumerate}[a)]
\item 设 $C$ 是带普通二重点 $x$ 的代数曲线，$C'$ 是其正规化，$a$ 和 $b$ 是 $C'$ 中位于 $x$ 之上的两个点. 设 $C'_i$ ($i = 1, 2$) 是 $C'$ 的两份拷贝，$a_i$ 和 $b_i$ 是 $C'_i$ 中分别对应于 $a$ 和 $b$ 的点. 在和曲线 $C'_1 \amalg C'_2$ 中，把 $a_1$ 与 $b_2$ 等同，把 $a_2$ 与 $b_1$ 等同 (把具体的等同过程留给读者；它将在大全第六章中说明，而在代数闭域上的曲线的情形，已在 Serre 关于代数曲线的书中处理). 得到一个连通但可约的曲线 $C''$，它是 $C$ 的 2 次平展覆盖. 读者将验证：一般地，$C$ 的连通``Galois 型''平展覆盖 $C''$，其原像 $C'' \times_C C'$ 是 $C'$ 的平凡覆盖 (i.e. 同构于若干份 $C'$ 的拷贝的和)，是 $n$ 次``循环''的，且对每个整数 $n > 0$，可以构造 $C$ 的一个 $n$ 次连通循环平展覆盖. 在稍后将发展的基本群语言中，这意味着 $\pi_1(C)$ 除以由 $\pi_1(C') \to \pi_1(C)$ (投影诱导的同态) 的像生成的闭不变子群所得的商，同构于 $\mathbf{Z}$ 的紧化 (compactification). 更精确地说，应当能够证明 $C$ 的基本群同构于 $C'$ 的基本群与 $Z$ 的紧化的 (拓扑) 自由积 (free product). 注意，正是这类问题催生了概形的``下降理论''.
\item 设 $A$ 是整完备局部环，我们知道其正规化 $A'$ 在 $A$ 上有限 (Nagata)，故它是完备半局部环，又因它整故是局部的. 假设它所定义的剩余扩张 $L/k$ 非根式 (在相反的情形，称 $A$ 几何单支，参见下文). 例如环 $\mathbf{R}[\![s, t]\!]/(s^2 + t^2)\mathbf{R}[\![s, t]\!]$ 就是如此，其中 $\mathbf{R}$ 是实数域. 设 $k'$ 是 $k$ 的有限 伽罗瓦扩张，使 $L \otimes_k k'$ 分解；设 $B$ 是对应于剩余扩张 $k'$ 的 $A$ 上的有限平展代数 (回忆 $B$ 本质上唯一). 则 $B$ 上的 $B' = A' \otimes_A B$ 有剩余代数 $L \otimes_k k'$，它不是局部的，故 $B'$ 不是局部环，从而 (完备) 有零因子 (zero divisors). 由于它包含在 $B$ 的全分式环 (total ring of fractions) 中 (因为它在 $A'$ 上自由，故在 $A'$ 上无挠 (torsion-free)，故在 $A$ 上无挠，故包含在 $B' \otimes_A K = B'_{(K)} = A'_{(K)} \otimes_K B_{(K)} = B_{(K)}$ 中，因为 $A'_{(K)} = K$)，由此推出 $B$ 不整. 对环 $\mathbf{R}[s, t]/(s^2 + t^2)\mathbf{R}[s, t]$，取 $k'/k = \mathbf{C}/\mathbf{R}$，得到的 $B$ 是平面中两条相交直线在交点处的局部环.
\end{enumerate}

还注意到，若存在 $Y$ (整) 的连通平展覆盖 $X$ 不是不可约的，则 $X$ 的每个不可约分支都给出支配 $Y$ 的、在 $Y$ 上非分歧但不在 $Y$ 上平展的覆盖 $X'$ 的例子. 在例子 a) 的情形，由此得到 $C'$ 在 $C$ 上非分歧，但在 $a$ 和 $b$ 两点处不平展 (这也可以直接通过检查 $x$ 和 $a$ 的局部环的完备化看出：从``形式''的观点，$C'$ 在 $a$ 点等同于 $C$ 在 $x$ 点的闭子概形，即经过 $x$ 的 $C$ 的两条``分支''之一).

在 a) 和 b) 中，我们看到 \ref{thm:I.9.5} (i) 和 (ii) 的结论不成立，这直接与``$Y$ 的一个点'爆裂'成正规化中的不同点''这一事实相关 (在 b) 中，剩余扩张非根式这一事实应当从几何上作如此解释). 精确地说，称整局部环 $A$ 几何单支 (geometrically unibranched)，若其正规化只有一个极大理想，且相应的剩余扩张是根式的；整预概形的点 $y$ 称为几何单支的，若其局部环几何单支. 例子：正规点，曲线的普通尖点 (cusp)，等等. 似乎若 $Y$ 有一个非单支的点，则 $Y$ 总存在一个连通非不可约的平展覆盖；至少在 $Y$ 是完备局部环的谱的情形，这就是我们在 b) 中所证明的. 反之可以证明，若 $Y$ 的所有点都几何单支，则每个在 $Y$ 上非分歧且支配 $Y$ 的连通 $Y$-预概形都平展且不可约. 证明沿用 \ref{thm:I.9.5} 的证明，使用下述定理 \ref{thm:I.8.3} 的推广，它将在稍后用下降技术证明(...)

设 $Y' \to Y$ 是有限、根式、满射的态射 (i.e. 可以称之为``泛同胚 (universal homeomorphism)''). 考虑从 $Y$-预概形到 $Y'$-预概形的函子 $X \mapsto X \times_Y Y' = X'$. 该函子诱导平展 $Y$-概形的范畴与平展 $Y'$-概形的范畴之间的等价. 例如，可以把这一结果应用于 $Y'$ 是 $Y$ 的正规化、$Y$ 假设为单支 (且 $Y'$ 在 $Y$ 上有限，这在实践中所遇到的所有情形都成立) 的情形，或应用于``夹在'' $Y$ 与其正规化之间的 $Y''$ (它不再需要在 $Y$ 上有限) 的情形.


